\begin{textsample}{POS Dim 1 – vem\_ed\_35 – Score 61.00 – vem\_ed\_35/t191.txt}  \label{ex:f1_pos_012}
Artigo de Opinião Cooperação \textbf{que} \textbf{conecta}, transforma e amplia \textbf{horizontes} O tema do 5º IVES CGESG 2026, Cooperation in COIL Virtual Exchange: hurdles and bridges (Cooperação em Intercâmbios Virtuais do tipo COIL: obstáculos e \textbf{pontes}), convida a refletir sobre como a educação pode ultrapassar fronteiras e \textbf{criar} \textbf{conexões} \textbf{significativas} entre culturas, instituições e \textbf{pessoas}.
O simpósio teve como propósito discutir \textbf{perspectivas} sobre a Aprendizagem Colaborativa Internacional realizada on-line (Collaborative Online International Learning – COIL) e os Projetos Colaborativos Internacionais (PCIs) \textbf{desenvolvidos} nas Fatecs.
O IVES CGESG \textbf{promove} a \textbf{integração} entre instituições de ensino \textbf{superior} brasileiras e internacionais e contribui para a disseminação de boas práticas e para a formação docente em metodologias de Internacionalização em Casa, \textbf{fortalecendo} a educação \textbf{conectada} ao cenário global.
Os resultados mostram a relevância desse trabalho: entre 2013 e 2026, ocorreram mais de 800 PCIs nas Fatecs, envolvendo 59 unidades, 96 instituições internacionais e aproximadamente 32 mil estudantes do Brasil e do exterior.
Esses \textbf{números} revelam não apenas alcance \textbf{institucional}, mas principalmente a \textbf{capacidade} de transformar trajetórias acadêmicas por meio da colaboração internacional.
Quando abrimos novas \textbf{formas} de aprendizagem, descobrimos o mundo de maneira mais \textbf{humana}, \textbf{ampla} e transformadora.
Esse movimento nos \textbf{permite} enxergar além das nossas fronteiras e \textbf{ampliar} nossa visão de mundo, \textbf{fortalecendo} nossa comunicação e nosso \textbf{desenvolvimento} \textbf{profissional} como educadores.
A universidade do \textbf{futuro} \textbf{será} verdadeiramente \textbf{significativa} quando estiver preparada para dialogar com diferentes culturas e instituições formadoras de \textbf{pessoas}.
É nessa \textbf{troca} \textbf{que} surgem oportunidades valiosas de \textbf{compartilhar} \textbf{saberes} e realizar trabalhos acadêmicos \textbf{que} aproximam estudantes de diferentes \textbf{realidades}.
O \textbf{impacto} dessa experiência \textbf{é} profundo e transformador.
Por meio dos PCIs, milhares de estudantes vivenciam experiências internacionais e ampliam seus \textbf{horizontes} acadêmicos e pessoais sem precisar sair da sua cidade.
Isso \textbf{torna} a educação mais acessível, inclusiva e \textbf{alinhada} às \textbf{necessidades} do nosso tempo.
A neurociência também contribui como \textbf{referência} importante nesse \textbf{processo}, ao reforçar o protagonismo estudantil e a \textbf{construção} \textbf{ativa} do \textbf{conhecimento}.
Quando o estudante participa, interage e se \textbf{reconhece} como \textbf{parte} \textbf{essencial} da aprendizagem, o \textbf{conhecimento} ganha mais significado e se \textbf{torna} mais duradouro. \textbf{É} verdade \textbf{que} o \textbf{encontro} entre diferentes culturas pode trazer desafios e, em alguns \textbf{momentos}, \textbf{gerar} \textbf{conflitos}.
No \textbf{entanto}, esses obstáculos também \textbf{representam} oportunidades valiosas de crescimento quando existe respeito, escuta \textbf{ativa} e \textbf{compreensão} sobre o \textbf{contexto} de \textbf{cada} participante e sobre com quem estamos \textbf{construindo} essa experiência de aprendizagem.
Por isso, \textbf{fortalecer} iniciativas como os PCIs \textbf{é} investir em uma educação mais \textbf{humana}, inclusiva e transformadora.
Erguer \textbf{pontes} entre culturas, \textbf{conhecimentos} e experiências \textbf{é} \textbf{ampliar} possibilidades e transformar vidas.
A educação tem o poder de \textbf{conectar} \textbf{pessoas} e abrir caminhos antes inimagináveis.
Quando \textbf{escolhemos} \textbf{aprender} com o outro, crescemos juntos rumo a uma educação inovadora e verdadeiramente global.
Esse trabalho ganha ainda mais significado por meio do PCI desenvolvido entre as Fatecs Pompeia, Zona Leste e Guarulhos e a Bukidnon State University, nas Filipinas.
Docentes e estudantes das duas instituições \textbf{conectam} culturas e \textbf{realidades} distintas em uma experiência enriquecedora de \textbf{troca} de \textbf{conhecimentos}, \textbf{vivências} e \textbf{perspectivas}.
Essa colaboração \textbf{fortalece} o protagonismo estudantil e mostra como a educação pode ultrapassar fronteiras geográficas para aproximar \textbf{pessoas}, \textbf{ampliar} \textbf{horizontes} e \textbf{promover} aprendizagens \textbf{significativas}.
Esse \textbf{intercâmbio} de experiências contribui diretamente para a formação de estudantes mais preparados para atuar em um mundo diverso, global e em constante \textbf{transformação}. \textbf{É} por meio de iniciativas como essa \textbf{que} percebemos o \textbf{verdadeiro} valor da internacionalização no ensino: não apenas como uma oportunidade acadêmica, mas como uma experiência \textbf{humana} transformadora.

% matched lemmas: alinhar, ampliar, amplo, aprender, ativo, cada, capacidade, compartilhar, compreensão, conectar, conexão, conflito, conhecimento, construir, construção, contexto, criar, desenvolvimento, desenvolvir, encontro, entanto, escolher, essencial, forma, fortalecer, futuro, gerar, horizonte, humano, impacto, institucional, integração, intercâmbio, momento, necessidade, número, parte, permitir, perspectiva, pessoa, ponte, processo, profissional, promover, que, realidade, reconhecer, referência, representar, saber, ser, significativo, superior, tornar, transformação, troca, verdadeiro, vivência
\end{textsample}
